% !TeX root = ../document_maitre.tex

\label{uxui}


Avant toute choses, notons qu'il s'agit ici d'une partie réflexive qui n'a pas pour objectif d'aboutir à des solutions concrètes. L'ambition est de présenter un état des lieux et des pistes de réflexions pouvant être l'objet de logiques d'implémentations et d'usage de ces outils.\newline

Lors de notre exploration du \gls{chatbot}, deux notions fondamentales ont été évoquées : l'\gls{ux} et l'\gls{ui}. Ils visent à placer l'utilisateur au c\oe{ur} des conceptions logicielles car l'usage utilisateur offerte défini l'attractivité et la qualité logicielle. Ils visent alors à garantir une utilisation optimale, fluide et intuitive d'un logiciel par le respect d'un certains nombres de principes, de normes et de méthodes. On y distingue l'expérience, gérée par l'\gls{ux}, et l'interface, gérée par l'\gls{ui}.

Arriver à garantir ces aspects permet de renforcer l'adoption du logiciel au près du public cible. Au c\oe{ur} de ces aspects se trouve l'expérience utilisateur, c'est à dire l'ensemble de ressenti, sensations et \guillemet{émotions} qu'une utilisation de l'application produit ou doit produire. L'\gls{ux} vise à garantir une expérience globale positive et qualitative lors de l'interaction entre l'utilisateur et l'application.\footcite{groupeLUXDesignDefinition2022} Un certains nombre de principes et de règles\footnote{Notamment \cite{ISO92411122025}} constituant ce domaine sont déterminées dans cet objectif. L'\gls{ux} régit un aspect pratique créé au moment de l'élaboration et du codage du logiciel, mobilisant ainsi un certains nombre de méthodes pour concevoir une expérience idéale dans une application orientée utilisateur.\footcite{lallemandMethodesDesignUX2018}

L'expérience utilisateur est transmise au travers d'un intermédiaire spécifique qu'est l'interface. Pour l'utilisateur l'interface fait partie intégrante de l'expérience globale, du côté des développeurs l'interface se différencie de l'expérience car n'étant qu'un outil de diffusion et d'interactivité. Ainsi, la création d'interface qualitative permettant une interaction \guillemet{positive} entre l'utilisateur le logiciel est régi par son domaine propre : l'\gls{ui}. Les fondements de ce dernier sont profonds. Dès 1990, le psychologue Don Norman démontre la responsabilité des \textit{design} dans les erreurs humaines dans un contexte industriel : \textit{\guillemet{Most industrial accidents are caused by human error: estimates range between 75 and 95 percent. How is it that so many people are so incompetent? Answer: They aren’t. It’s a design problem.}}\footcite{normanDesignEverydayThings} C'est le fondement de ce domaine : une interface mal conçue entraîne nécessairement une erreur. L'\gls{ui} a pour objectif de traduire l'\gls{ux} via des repères visuels et des parcours graphiques permettant d'activer des fonctionnalités. À l'instar de la notion précédente, l'\gls{ui} s'appuie sur des principes graphiques pour créer des interfaces intuitives au service de l'utilisateur. \newline

Soulignons que ces domaines ne sont pas rigide et ne posent pas de contraintes ni d'obligations. Ils ne font que proposer des méthodologies et outils obéissant à des principes et règles pensées pour atteindre des objectifs. Chaque développement logiciel n'a pas d'obligation à les utiliser ou à en respecter les principes. Ils peuvent adapter ces éléments où bien créer leurs propres principes, règles et méthodologies. 

Prenons l'exemple de la norme ISO 9241-112 abordant la question de l'expérience et interface des systèmes informatiques. Elle ne fixe pas de cadre rigide et obligatoire, elle \guillemet{[...] énonce des principes de conception ergonomique pour des systèmes interactifs dont la présentation d'information de l’interface utilisateur est contrôlée par logiciel}\footcite{ISO92411122025}. Son objectif est de proposer des principes clés pour conceptualiser une structure logicielle optimisée pour la compréhension des informations par l'utilisateurs. Il subsiste une certaine liberté et un flou conceptuel qui rendent l'application de ces concepts parfois délicate.

\subsection{\gls{ux} et \gls{ui} dans les \glspl{sgc} :}

Nous l'avons déjà exploré, les \glspl{sgc} sont des applications aux architectures denses et complexe ce qui impact nécessairement l'expérience et l'interface. Ils sont basés sur l'intégration d'une multiplicité d'informations avec une haute interactivité entre elles réparties dans de multiples fonctionnalités. Une notice d'objet disposera de ces propres fonctionnalités et interfaces, de même pour une notice de prêt. Ainsi, les interfaces sont multiples car elles correspondent à des expériences multiples adaptées à des modalités de gestion spécifiques. Cette multiplicité de parcours entraîne des complexités structurelles dans l'organisation de l'application et des données contenues et diffusées.

Dans ce contexte, l'\gls{ux} et l'\gls{ui} sont fondamentales pour simplifier l'apparence et l'usage d'outils destinés à manipuler de très grands volumes de données et à appuyer des décisions stratégiques. Il n'existe aucun principes, règles ou méthodologies spécifiques à l'application de ces domaines aux \glspl{sgc}. Cela s'explique par l'existence de questionnement et de contradiction, entre les principes généraux de ces domaines et les besoins auxquels répondre. La conception de ces domaines vise à simplifier l'usage applicatif et l'affichage par une logique d'épuration pour plus d'efficacité. Cependant, les besoins patrimoniaux au c\oe{ur} de la raison d'être de ces logiciel est l'exhaustivité des données ou plutôt la possibilité d'interagir avec des données dans toutes leurs complexités. Cette logique amène à dupliquer les expériences utilisateurs et à charger les interfaces avec deux informations : des fonctionnalités, représentées par des boutons dit d'actions, et les données stockées. 

\textbf{inclure une photo}\newline 

Aujourd'hui, des discussions académiques et professionnelles amènent à repenser ces notions appliquées aux \gls{sgc}. Deux idées principales en émergent : l'application du concept de simplicité sans tendre vers l'épuration, limitant l'exhaustivité naturelles de ces collections, et la recherche de nouvelles solutions d'interactions et graphiques. L'idée étant de simplifier l'exhaustivité, et non de la limiter et de la biaiser ce qui est aujourd'hui le cas lorsque l'on parle de rationalisation des interfaces et expériences utilisateurs. 

De plus, on ne s'adresse pas à des utilisateurs grands publics mais à des professionnels aux objectif précis et aux contraintes professionnelles devant être prise en compte. Ces objectifs et contraintes professionnelles varient en fonction de la typologie de collections traitées et de l'aspect de la collection qui est traité. On peut imaginer qu'une collection d'histoire naturelle ne demandera pas les mêmes expériences et interface que pour une bibliothèques. \newline

Aujourd'hui, des tentatives de compromis entre simplicité, exhaustivité et particularité de chaque collections tendant à apparaître dans une logique de communication vers un large public.\footcite{whitelawGenerousInterfacesDigital2015} On peut tout à fait s'attendre à voir émerger le souhait d'une cohérence de traitement, poussant les \gls{sgc} dans cette direction. A savoir : l'utilisation de moyens graphiques \textit{(comme des graphes)} pour représenter l'envergure de la collection et s'adapter à des usages et objectifs précis.\footcite{whitelawGenerousInterfacesDigital2015}

Aujourd'hui, les ambitions d'\gls{ux} et d'\gls{ui} appliquées à des \gls{sgc} semblent suivre la voix tracée par deux éléments. D'abord, la théorie de la surcharge cognitive\footcite{paasCognitiveLoadTheory2010} qui démontre qu'une surcharge informationnelle impact directement la productivité utilisateurs en réduisant les capacité d'attentions et de concentration de l'utilisateur. Transposé à un environnement logiciel, un surplus d'outils et de données risque d'entraîner une saturation chez l'utilisateur conduisant alors a un rejet du système. Ce phénomène est au cœur du (\textit{Technology Acceptance Model})\footcite{baddaTechnologyAcceptanceModel2025}, rappelant que l'utilité perçue et la facilité d'utilisation conditionnent directement l'adoption d'un outil. Modèle vers lequel les logiciels présentant de tels densités d'informations tendent aujourd'hui à appliquer. 


\subsection{Les outils d'\gls{IA} :}

L'inclusion de l'\gls{IA} dans toutes sortes de logiciel bouscule la réalité de l'\gls{ux} et de l'\gls{ui}. En effet, il est tentant de considérer le développement de fonctionnalités basées sur l'\gls{IA} comme identique aux autres développement mais ce n'est pas tout à fait la réalité. \newline

Les outils d'\gls{IA} sont souvent développés par des entreprises externes sans lien avec le système souhaitant l'inclure. Ils sont donc pensés indépendamment de systèmes tierces, avec leurs propres interfaces et expériences. Ceci pose un premier défis qu'est celui de l'inclusion de l'outil à inclure dans le système. L'outil possède ses propres besoins et propositions d'expérience et d'interface de même que le système hôte ce qui peux créer un important décalage entre les éléments et poussant à adapter son \gls{ux} et \gls{ui} d'origine afin de créer un dialogue entre les deux éléments. 

Ajoutons également que ces outils entendent parfois \guillemet{révolutionner} certaines pratiques et procédure communes. Ceci pose la question des habitudes utilisateurs et de leurs changements, ce qui peux impacter la qualité d'expérience globale. En effet, trop changer en profondeur une fonctionnalité habituelle peu déstabiliser un utilisateur, parfois définitivement. Par exemple, l'apparition d'un \gls{chatbot} peut être destabilisant : on est dès lors capable de converser à travers une fenêtre dédiée et de l'utiliser pour effectuer des tâches différemment, voir parfois les automatiser.

Pour finir, l'utilisation d'\gls{IA} pose aussi la question de la transparence. Nous l'avons vu, ces outils agissent sur des données parfois personnelles et/ou sensibles et des questions légales et éthiques impose une obligation de transparence. Cela dans le but que l'utilisateur puisse savoir les données utilisées, dans quels buts ses données sont utilisées notamment. Cela demande donc de créer une expérience rassurante et transparente, amenant à modifier les interfaces pour permettre la diffusion de ces nouvelles informations.\newline

Ces quelques points nous permettent de comprendre que l'intégration d'\gls{IA} dans des logiciels demande plus que la création d'une fonctionnalité technique. Il faut aussi, et surtout, réussir à intégrer de nouvelles expériences et interfaces rassurantes dans des cadres préexistants sans perturber l'utilisateur. C'est une tâche difficile dont dépend l'adoption de ces outils par les utilisateurs. 

\subsection{\gls{IA} et \gls{sgc} : des attentes \gls{ux} et \gls{ui} particulières}

L'inclusion d'outils basés sur l'\gls{IA} dans des \gls{sgc} pose des problématiques particulières. Les utilisateurs sont des professionnels souvent alerte sur l'utilisation des données par ce type d'outils, sensibles à la conservation de leurs données et engagés dans des démarches de rigueur scientifique. Ainsi, trois aspects semblent être au c\oe{ur} de tensions : \newline

D'abord la confiance envers l'outils. Malgré des efforts de contrôles, l'hallucination reste fréquentes et gênantes. La gestion des données patrimoniales est gouvernée par un besoin d'exactitude scientifique autant dans sa production que dans sa recherche. L'hallucination peu compromettre cette rigueur et entraîner une méfiance, un désintérêt légitime voire un rejet de l'outil. L'\gls{ux} doit pouvoir prendre en compte ses éléments en proposant des solutions comme des processus dit \textit{humain in the loop} et de gestion des erreurs, plaçant ainsi l'utilisateur comme contrôleur de la qualité des données et donc en position de contrôle. 

Le second élément est celui de l'interactions de ces outils avec l'environnement logiciel préexistant. Par exemple, la recherche augmentée comme le \gls{chatbot} génèrent de la donnée en masse pouvant amener à créer une surcharge informationnelles. La problématique se situe dans l'insertion de ces outils dans des environnements métiers déjà maîtrisé par l'utilisateur. Il faut s'assurer que ces outils amènent une plus values, qu'ils permettent une simplification des parcours d'usages sans créer des incohérences de parcours.

Le dernier élément est celui de la perception. Notre examen de la recherche intelligente démontre que le public auxquels on s'adresse aujourd'hui est méfiant envers les \glspl{IA}. Ceci peut expliquer partiellement le manque de travaux sur l'utilisation de \gls{chatbot} dans la production et la gestion de données patrimoniales. Faire de ces outils de véritables assistants à la productivité plutôt que des contraintes visuelles et technique constitue un enjeu majeur.

\subsection{Conclusion}

L'intégration du \gls{NLP} en deux outils différents, la recherche augmentée et le \gls{chatbot}, soulève des problématiques d'\gls{ux} et d'\gls{ui} inédites. Il ne s'agit pas d'ajouter des fonctionnalités supplémentaires ou de modifier légèrement une expérience utilisateur et l'interface. Il s'agit de changer drastiquement des éléments préexistant par un changement de rapport et de dialogue entre l'utilisateur et son système professionnel dans un domaine exigeant et attentif aux questions de sécurités et de qualités. 

Ces nouveaux outils doivent donc amener la conception de nouvelles expériences et interfaces respectant les exigences professionnelles fortes et réduisant la surcharge visuelle et cognitive. Ou du moins, ne pas l'alourdir. 

Il faut penser à l'adoption de ces technologies plus qu'aux simples performances technologiques. Cette adoption passe par la nécessité d'un dialogue transparent entre l'outil et l'utilisateur, placé au c\oe{ur} du processus. Ce qui déplace les enjeux de l'\gls{ux} et l'\gls{ui} des questions visuelles et de confort vers des questions de confiances et d'accessibilités.